\section{问题三模型建立和求解}

\subsection{三弹联合时序模型}

设 FY1 的共享航向角和速度为 $\theta$ 与 $v$，第 $i$ 枚烟幕弹的投放时刻、引信延时和起爆时刻分别为 $t_{r,i}$、$\tau_i$ 与 $t_{e,i}=t_{r,i}+\tau_i$。用非负松弛量 $s_{12},s_{23}$ 编码投放间隔：
\begin{equation}
 t_{r,2}=t_{r,1}+1+s_{12},\qquad
 t_{r,3}=t_{r,2}+1+s_{23},\qquad s_{12},s_{23}\ge0.
\end{equation}
第 $i$ 枚烟幕弹的投放点和起爆点由同一条航迹确定：
\begin{align}
 \bm r_{r,i}&=\bm u_0+v\bm e(\theta)t_{r,i},\\
 \bm r_{e,i}&=\bm u_0+v\bm e(\theta)t_{e,i}
 +(0,0,-\tfrac12g\tau_i^2)^{\mathsf T}.
\end{align}
分别记三枚烟幕弹的有效集合为 $\Omega_i$。总时长必须取并集 $\Omega=\bigcup_i\Omega_i$ 的测度，不能直接累加单弹时长，否则重叠部分会被重复计算。

\subsection{优化模型汇总}

问题三的八维优化模型为
\begin{equation}
\begin{aligned}
&\max_{\bm y}\quad T_{\mathrm{eff}}^{(3)}(\bm y)
=\mu\!\left(\bigcup_{i=1}^{3}\Omega_i(\bm y)\right),\\[-1mm]
&\text{\rm s.t.}\quad
\left\{\begin{aligned}
&\bm y=(\theta,v,t_{r,1},s_{12},s_{23},\tau_1,\tau_2,\tau_3)^{\mathsf T},\\
&0\le\theta<2\pi,\qquad 70\le v\le140,\\
&t_{r,1}\ge0,\quad s_{12},s_{23}\ge0,\\
&t_{r,2}=t_{r,1}+1+s_{12},\quad t_{r,3}=t_{r,2}+1+s_{23},\\
&\tau_i\ge0,\qquad t_{e,i}=t_{r,i}+\tau_i\le t_M,\\
&\bm r_{r,i}=\bm u_0+v\bm e(\theta)t_{r,i},\\
&\bm r_{e,i}=\bm u_0+v\bm e(\theta)t_{e,i}+(0,0,-\tfrac12g\tau_i^2)^{\mathsf T},\\
&z_{e,i}=1800-\tfrac12g\tau_i^2\ge0,\\
&\bm c_i(t)=\bm r_{e,i}-(0,0,3(t-t_{e,i}))^{\mathsf T},\\
&t_{e,i}\le t\le\min\{t_{e,i}+20,t_M\},\\
&D_i(t)=\displaystyle\max_{\bm p\in\mathcal T}d_i(\bm p,t),\\
&\Omega_i=\{t:D_i(t)\le10\},\qquad i=1,2,3.
\end{aligned}\right.
\end{aligned}
\label{eq:q3-complete-model}
\end{equation}
该模型同时约束共享航迹、投放间隔、自由落体高度损失、烟幕寿命和完整圆柱遮蔽几何。由松弛量编码产生的候选解天然满足最小投放间隔，可减少惩罚函数造成的无效搜索。

\subsection{联合优化与精确复核}

真实目标由多个阈值事件和区间并运算构成，具有大面积零值区、多个航向盆地和窄时间窗。搜索阶段先用平滑遮蔽概率近似引导差分进化，再在全航向、外向和内向扇区执行多初值搜索；同时加入首尾时间窗拼接构造，避免种群遗漏靠近区间衔接边界的高质量策略。差分进化适合处理不可导、多峰连续优化问题\cite{storn1997de}。

候选排序使用真实几何判据：计算每枚烟幕球到完整圆柱各条视线段的最大距离，以 Brent 法求解 $D_i(t)=10$ 的进出事件\cite{brent1973algorithms}，再合并全部有效区间。最终以每个圆周 $5760$ 个角向点和 $0.005\,\mathrm{s}$ 的扫描步长复核。流程见\cref{fig:q3-flow}。

\begin{figure}[!htbp]
 \centering
 \includegraphics[width=0.72\textwidth]{fig_flow_q3.pdf}
 \caption{问题三三弹联合优化与区间并集复核流程}
 \label{fig:q3-flow}
\end{figure}

\subsection{投放策略与遮蔽结果}

最优候选取航向角 $179.647479^{\circ}$、速度 $139.998280\,\mathrm{m/s}$，对应速度向量为
\begin{equation}
 \bm v_U=(-139.995630,0.861355,0)^{\mathsf T}\,\mathrm{m/s}.
\end{equation}
三次投放分别发生在 $0.003$、$3.702$ 和 $5.569\,\mathrm{s}$，相邻间隔为 $3.699$ 和 $1.867\,\mathrm{s}$。具体策略见\cref{tab:q3-strategy}。

\begin{table}[!htbp]
 \centering
 \caption{问题三 FY1 三枚烟幕弹投放策略}
 \label{tab:q3-strategy}
 \resizebox{\textwidth}{!}{
 \begin{tabular}{cccccc}
  \toprule[1.5pt]
  编号 & 投放时刻/s & 投放点/m & 起爆时刻/s & 起爆点/m & 单弹时长/s \\
  \midrule[1pt]
  1 & $0.003$ & $(17799.580,0.003,1800.000)$ & $3.614$ & $(17294.056,3.113,1736.107)$ & $3.834877$ \\
  2 & $3.702$ & $(17281.736,3.189,1800.000)$ & $9.039$ & $(16534.579,7.786,1660.431)$ & $2.591030$ \\
  3 & $5.569$ & $(17020.364,4.797,1800.000)$ & $11.610$ & $(16174.651,10.000,1621.181)$ & $1.204687$ \\
  \bottomrule[1.5pt]
 \end{tabular}}
\end{table}

三枚弹的有效区间分别为
\begin{equation}
 [5.205021,9.039899],\quad[9.039000,11.630030],\quad[11.610000,12.814687].
\end{equation}
相邻区间在端点附近交叠，因而拼接成
\begin{equation}
 \Omega=[5.205021,12.814687],\qquad
 T_{\mathrm{eff}}^{(3)}=7.609665451\,\mathrm{s}\approx\boxed{7.610\,\mathrm{s}}.
\end{equation}
与早投边界候选相比，新方案通过延迟起爆把三个云团依次布置在 M1 视线前方，第三枚弹获得正边际收益并延长尾部窗口。该方案同步写入 \texttt{result1.xlsx}，并作为问题五中 FY1 对 M1 的子策略。

\begin{figure}[!htbp]
 \centering
 \includegraphics[width=0.96\textwidth]{q3_optimization_diagnostics.pdf}
 \caption{问题三多盆地搜索与三弹最坏视线距离}
 \label{fig:q3-diagnostics}
\end{figure}

\subsection{收敛性与稳健性检验}

角向采样从 $1440$ 点加密到 $5760$ 点时，联合时长变化小于 $10^{-4}\,\mathrm{s}$。三个起爆高度分别为 $1736.107$、$1660.431$ 和 $1621.181\,\mathrm{m}$，均远高于地面；全部有效区间都处于相应烟幕起爆后的 $20\,\mathrm{s}$ 内，并早于 M1 到达假目标。

最优区间的衔接裕量较小，说明时序误差比位置小扰动更敏感。本文通过事件求根而非固定网格计数确定端点，避免把端点接触误判为断裂。不同随机种子可得到参数略有差异但时长相近的候选，因此本文将其表述为经多盆地搜索和高密度复核的最优候选，不声称解析意义上的全局最优。

\endinput
\section{问题三模型建立和求解}

\subsection{三弹联合时序模型}

设 FY1 的共享航向角和速度仍为 $\theta$ 与 $v$，第 $i$ 枚烟幕弹的投放时刻、引信延时和起爆时刻分别为 $t_{r,i}$、$\tau_i$ 与 $t_{e,i}=t_{r,i}+\tau_i$。为在搜索过程中严格满足相邻投放间隔约束，引入非负松弛量 $s_{12},s_{23}$，令
\begin{equation}
    t_{r,2}=t_{r,1}+1+s_{12},\qquad
    t_{r,3}=t_{r,2}+1+s_{23},\qquad s_{12},s_{23}\ge0.
    \label{eq:q3-release-code}
\end{equation}
无人机的单位航向向量为 $\bm e(\theta)=(\cos\theta,\sin\theta,0)^{\mathsf T}$。第 $i$ 枚烟幕弹的投放点和起爆点由同一条航迹确定：
\begin{align}
    \bm r_{r,i}&=\bm u_0+v\bm e(\theta)t_{r,i},\\
    \bm r_{e,i}&=\bm u_0+v\bm e(\theta)t_{e,i}
    +(0,0,-\tfrac12g\tau_i^2)^{\mathsf T}.
    \label{eq:q3-release-explosion}
\end{align}
起爆后，第 $i$ 个烟幕球心以 $3\,\mathrm{m/s}$ 竖直下沉。其有效时间集合记为 $\Omega_i$；三弹联合遮蔽集合必须取并集 $\Omega=\bigcup_{i=1}^{3}\Omega_i$，从而避免重复计算重叠时间。

\subsection{优化模型汇总}

把共享航迹、投放间隔、自由落体、烟幕寿命、导弹到达时限和完整圆柱遮蔽判据统一写入约束，问题三的优化模型为
\begin{equation}
\begin{aligned}
&\max_{\bm y}\quad T_{\mathrm{eff}}^{(3)}(\bm y)
=\mu\!\left(\bigcup_{i=1}^{3}\Omega_i(\bm y)\right),\\[-1mm]
&\text{\rm s.t.}\quad
\left\{
\begin{aligned}
&\bm y=(\theta,v,t_{r,1},s_{12},s_{23},\tau_1,\tau_2,\tau_3)^{\mathsf T},\\
&0\le\theta<2\pi,\qquad 70\le v\le140,\\
&t_{r,1}\ge0,\quad s_{12},s_{23}\ge0,\\
&t_{r,2}=t_{r,1}+1+s_{12},\quad
 t_{r,3}=t_{r,2}+1+s_{23},\\
&\tau_i\ge0,\qquad t_{e,i}=t_{r,i}+\tau_i\le t_M,\quad i=1,2,3,\\
&\bm u(t)=\bm u_0+v\bm e(\theta)t,\qquad
 \bm r_{r,i}=\bm u(t_{r,i}),\\
&\bm r_{e,i}=\bm u_0+v\bm e(\theta)t_{e,i}
 +(0,0,-\tfrac12g\tau_i^2)^{\mathsf T},\\
&z_{e,i}=1800-\tfrac12g\tau_i^2\ge0,\qquad
 t_M=\dfrac{\lVert\bm m_0\rVert}{300},\\
&\bm c_i(t)=\bm r_{e,i}-(0,0,3(t-t_{e,i}))^{\mathsf T},\\
&t_{e,i}\le t\le\min\{t_{e,i}+20,t_M\},\\
&\mathcal T=\{(x,y,z):x^2+(y-200)^2\le49,\ 0\le z\le10\},\\
&D_i(t;\bm y)=\displaystyle\max_{\bm p\in\mathcal T}d_i(\bm p,t;\bm y),\\
&\Omega_i(\bm y)=\{t:D_i(t;\bm y)\le10,\quad
 t_{e,i}\le t\le\min(t_{e,i}+20,t_M)\}.
\end{aligned}
\right.
\end{aligned}
\label{eq:q3-complete-model}
\end{equation}

式\cref{eq:q3-complete-model} 的目标是三个有效集合并集的测度，而非三个单弹时长之和。主目标相同的候选方案按两级辅助准则筛选：先优先选择 $\Omega_1$ 与 $\Omega_2$ 重叠裕量较大的方案，以降低小幅时序误差造成区间断裂的风险；若仍相同，则最小化末次投放时刻 $t_{r,3}$，使任务尽早完成。这些辅助准则只用于等时长平台内选取代表解，不改变主目标。

\subsection{联合优化与精确复核}

三弹决策空间为八维，真实目标由多个阈值事件和区间并运算构成，存在大面积零值区和投放下界附近的窄盆地。搜索阶段以平滑化的遮蔽概率并集作为引导，在全航向、外向扇区和内向扇区分别执行固定种子的差分进化，并用 Powell 法细化候选；同时显式加入 $t_{r,1}=0$、$t_{r,2}=1$、$t_{r,3}=2$ 的时序边界构造解，防止群体搜索遗漏活跃边界。

候选排序仍使用真实几何判据：对真目标上下圆周加密采样，计算每枚烟幕球到所有有限视线段的最大距离，以 Brent 法求解 $D_i(t)=10$ 的进出事件，再合并三个有效区间。最终以每个圆周 $5760$ 个角向点和 $0.005\,\mathrm{s}$ 的事件扫描步长复核，并作局部扰动、去弹边际贡献和等时长代表方案对比。求解流程见\cref{fig:q3-flow}。

\begin{figure}[!htbp]
    \centering
    \includegraphics[width=0.72\textwidth]{fig_flow_q3.pdf}
    \caption{问题三三弹联合优化与区间并集复核流程}
    \label{fig:q3-flow}
\end{figure}

\subsection{投放策略与遮蔽结果}

在联合时长相同的参数平台内，选取航向角 $\theta=0.097738438\,\mathrm{rad}=5.600000^{\circ}$、速度 $v=125.000000\,\mathrm{m/s}$ 的稳健代表方案。FY1 速度向量为
\begin{equation}
    \bm v_U=(124.403425,\ 12.197862,\ 0)^{\mathsf T}\ \mathrm{m/s}.
\end{equation}
三枚烟幕弹均采用即时起爆，投放时刻依次为 $0$、$1$、$2\,\mathrm{s}$。具体坐标与单弹有效时长见\cref{tab:q3-strategy}。

\begin{table}[!htbp]
    \centering
    \caption{问题三 FY1 三枚烟幕弹投放策略}
    \label{tab:q3-strategy}
    \resizebox{\textwidth}{!}{
    \begin{tabular}{cccccc}
        \toprule[1.5pt]
        编号 & 投放时刻/s & 投放点/m & 起爆时刻/s & 起爆点/m & 单弹有效时长/s \\
        \midrule[1pt]
        1 & $0.000$ & $(17800.000,0.000,1800.000)$ & $0.000$ & $(17800.000,0.000,1800.000)$ & $2.556003$ \\
        2 & $1.000$ & $(17924.403,12.198,1800.000)$ & $1.000$ & $(17924.403,12.198,1800.000)$ & $4.548404$ \\
        3 & $2.000$ & $(18048.807,24.396,1800.000)$ & $2.000$ & $(18048.807,24.396,1800.000)$ & $0.000000$ \\
        \bottomrule[1.5pt]
    \end{tabular}}
\end{table}

第 2 枚烟幕弹形成区间 $[1.000000,5.548404]$，第 1 枚形成区间 $[4.847157,7.403159]$，二者重叠 $0.701247\,\mathrm{s}$，故并集连续覆盖
\begin{equation}
    \Omega=[1.000000,7.403159],\qquad
    T_{\mathrm{eff}}^{(3)}=6.403159406\,\mathrm{s}
    \approx\boxed{6.403\,\mathrm{s}}.
\end{equation}
第 3 枚烟幕弹在该共享航迹上未形成完整圆柱遮蔽区间，对本代表方案的边际贡献为零；仍按题意完成第三次投放，并由末次投放时刻最小的辅助准则取 $t_{r,3}=2\,\mathrm{s}$。这不表示第三枚弹在所有可行航迹中均无作用，而是说明在当前单航向、单速度与完整目标判据下，前两枚弹已经构成本文搜索到的最长连续屏障。

\begin{figure}[!htbp]
    \centering
    \includegraphics[width=0.96\textwidth]{q3_optimization_diagnostics.pdf}
    \caption{问题三多盆地搜索与三弹最坏视线距离}
    \label{fig:q3-diagnostics}
\end{figure}

\subsection{非唯一性与稳健性检验}

加密计算表明，圆周采样从 $90$ 增至 $5760$ 时联合时长均保持 $6.403159406\,\mathrm{s}$。航向 $5.15^{\circ}$、速度 $140\,\mathrm{m/s}$ 的方案也达到相同主目标，但前两枚弹仅重叠约 $0.371008\,\mathrm{s}$；本文代表方案的重叠裕量增至 $0.701247\,\mathrm{s}$。因此不同算法可能给出航向角和速度不同、有效时长近乎一致的结果，其根源是目标函数存在等值平台，而不必然意味着计算错误。

局部扰动显示，将第 1 枚投放推迟 $0.01\,\mathrm{s}$ 时联合时长降为 $6.388954\,\mathrm{s}$，将第 2 枚投放推迟 $0.01\,\mathrm{s}$ 时降为 $6.393159\,\mathrm{s}$；负向扰动则违反投放时刻或最小间隔约束。三枚起爆点均位于地面以上，相邻投放间隔均恰为 $1\,\mathrm{s}$，所有有效区间均早于导弹到达假目标且处于烟幕寿命内。由于采用启发式全局搜索与边界筛查，本文将该结果表述为经多盆地和高密度复核的最优候选，而不宣称已经给出解析意义下的全局最优性证明。

\endinput
\section{问题三模型建立和求解}

\subsection{脚注}

用 \verb|\footnote{具体内容}| 可生成脚注\footnote{脚注可以补充说明一些东西}。

\subsection{无序列表与有序列表}

无序列表：
\begin{itemize}
  \item xelatex
  \item pdflatex
  \item lualatex
  \item etc.
\end{itemize}

有序列表：
\begin{enumerate}
  \item xelatex
  \item pdflatex
  \item lualatex
  \item etc.
\end{enumerate}

\subsection{字体加粗与斜体}

若想强调部分内容，可用加粗来实现。加粗字体可用 \verb|\textbf{加粗}| 来实现。如： \textbf{这是加粗的字体。 This is bold fonts} 。

中文字体没有斜体设计，但英文字体有。\textit{斜体 Italics}。
